%%%%%%%%%%%%%%%%%%%%%%%%%%%%%%%%%%%%%%%%%%%%%%%%%%%%%%%%%%%%%%%%%%%%%%
% Lớp 10 | Véctơ (chưa xét tọa độ) | Khái niệm véctơ
% Dạng 0H5?1-3: Hai véctơ bằng nhau
% Số câu: 9  (sinh tự động bởi tach_cau.py; từ nay thầy sửa trực tiếp file này)
%%%%%%%%%%%%%%%%%%%%%%%%%%%%%%%%%%%%%%%%%%%%%%%%%%%%%%%%%%%%%%%%%%%%%%

%%%%%%%%%%%%%%%% TRẮC NGHIỆM 4 PHƯƠNG ÁN %%%%%%%%%%%%%%%%

% nguồn: MaHoaNH_TracNghiem_Lop10
\begin{ex}%[0H5N1-3]%[ma:c34ce595]
	Khẳng  định nào sau đây đúng?
	\choice 
		{Hai vectơ $\overrightarrow{a}$ và $\overrightarrow{b}$ được gọi là bằng nhau nếu  chúng cùng phương}
		{Hai vectơ $\overrightarrow{a}$  và $\overrightarrow{b}$ được gọi là bằng nhau nếu chúng ngược hướng và cùng độ dài}
		{Hai vectơ $\overrightarrow{a}$ và $\overrightarrow{b}$ được gọi là bằng nhau nếu chúng cùng độ  dài}
		{\True Hai vectơ $\overrightarrow{a}$ và $\overrightarrow{b}$ được gọi là bằng nhau nếu  chúng cùng hướng và cùng độ dài}
	 \loigiai{
		Hai vectơ $\overrightarrow{a}$ và $\overrightarrow{b}$ được gọi là bằng  nhau nếu chúng cùng hướng và cùng độ dài.
	} 
\end{ex}

% nguồn: MaHoaNH_TracNghiem_Lop10
\begin{ex}%[0H5N1-3]%[ma:92daa729]
	Khẳng  định nào sau đây đúng?
	 \choice
		{\True Hai vectơ bằng nhau thì có giá trùng hoặc song  song}
		 {Hai vectơ có độ dài không bằng nhau thì không cùng hướng}
		{Hai vectơ không  bằng nhau thì chúng không cùng hướng}
		{Hai vectơ không bằng nhau thì độ dài của chúng không bằng nhau} 
	\loigiai{ 
		 Hai vectơ bằng nhau thì cùng hướng nên giá của chúng song song hoặc trùng nhau.
	} 
\end{ex}

% nguồn: MaHoaNH_TracNghiem_Lop10
\begin{ex}%[0H5N1-3]%[ma:906b13ed]
	 Cho lục giác đều $ABCDEF$ tâm $O$. Vectơ bằng vectơ $\overrightarrow{AB}$ là
	\choice 
		{$\overrightarrow{DE}$} 
		{$\overrightarrow{CO}$} 
		{\True  $\overrightarrow{FO}$}
		 {$\overrightarrow{AC}$}
	 \loigiai{
	\immini 
		 {
		Do $ABCDEF$ là lục giác đều  nên $FO\parallel AB$ và $FO=AB$. \\
		Suy  ra $\overrightarrow{FO} = \overrightarrow{AB}$.
		} 
		 {
		\begin{tikzpicture}[line join=round,line cap = round,>=stealth,scale=.8,font=\footnotesize,transform shape]
				 \foreach \x/\z/\g in
					{ 
					 A/180/180,B/120/135,C/60/45,D/0/0,E/-60/-45,F/-120/-135
				 }
				\draw[fill=black] ($(0,0)+(\z:3)$) circle (1pt) coordinate  (\x)
					 ($(\x)+(\g:3.5mm)$) node{$\x$};
				\draw[fill = black] (0,0) circle (1pt)  node[below]{$O$};
				 \draw
					(A)--(B)--(C)--(D)--(E)--(F)--(A) 
					(A)--(D) 
					 (B)--(E)
					(C)--(F); 
			\end{tikzpicture}
		 }
	} 
\end{ex}


%%%%%%%%%%%%%%%% ĐÚNG/SAI %%%%%%%%%%%%%%%%

% nguồn: MaHoaNH_DungSai_Lop10
\begin{ex}%[0H5H1-3]%[ma:22cf6b52]
	Cho tam giác $ABC$ có $G$ là trọng tâm. Gọi $G'$ là điểm đối xứng với $G$ qua trung điểm $M$  của $BC$.
	 \choiceTF
		{\True $\overrightarrow{GA}$, $\overrightarrow{G'G}$  cùng hướng}
		 {$GA=3GM$}
		 {\True $\overrightarrow{GA}=\overrightarrow{G'G}$}
		{\True $\overrightarrow{BG}=\overrightarrow{G'C}$, $\overrightarrow{BG'}=\overrightarrow{GC}$, $\overrightarrow{BM}=\overrightarrow{MC}$, $\overrightarrow{GM}=\overrightarrow{MG'}$, $\overrightarrow{AG}=\overrightarrow{GG'}$} 
	\loigiai{ 
		\begin{center}
			\begin{tikzpicture}[declare function={a = 5; goc = 65; b = 0.65*a;}]
					 \path
						 (goc:b) coordinate (A)
						(0,0) coordinate (B) 
						(a,0)  coordinate (C)
						($(B)!.5!(C)$) coordinate  (M)
						($(A)!2/3!(M)$)  coordinate (G)
						($(M)!-1!(G)$) coordinate  (G');
					\draw (A)--(B)--(C)--cycle  (A)--(G');
					\foreach \t/\g in  {A/90,B/-90,C/-90,M/220,G/0,G'/0}{
						\draw[fill=black]  (\t) circle (1pt) node[shift={(\g:7pt)},font=\scriptsize]{$ \t $};
					} 
				\end{tikzpicture}
		\end{center}
		\begin{itemchoice}
			\itemch Ta  có $\overrightarrow{GA}$, $\overrightarrow{G'G}$ cùng hướng.\hfill (1)
			\itemch Do $G$ là trọng tâm của tam giác $ABC$  nên $GA=2GM$.
			\itemch Ta có $G'$ là điểm đối xứng  với $G$ qua $M$ nên $G'G=2GM$. \hfill (2)\\
			Từ (1) và (2) suy ra  $\overrightarrow{GA}=\overrightarrow{G'G}$.
			\itemch $\overrightarrow{BG}=\overrightarrow{G' C}$, $\overrightarrow{BG'}=\overrightarrow{GC}$, $\overrightarrow{BM}=\overrightarrow{MC}$, $\overrightarrow{GM}=\overrightarrow{MG'}$,  $\overrightarrow{AG}=\overrightarrow{GG'}$.
		\end{itemchoice}
	 }
\end{ex}

% nguồn: MaHoaNH_DungSai_Lop10
\begin{ex}%[0H5H1-3]%[ma:eb6c470b]
	Cho tam giác $ABC$ có $M$, $N$ lần lượt là trung  điểm của $AB$ và $AC$. Lấy $P$ đối xứng với $M$ qua $N$.
	 \choiceTF
		{$MN = BC$} 
		 {\True $\left| \overrightarrow{MP} \right| = \left| \overrightarrow{BC} \right|$}
		 {$\overrightarrow{MN}$ và $\overrightarrow{BC}$ ngược hướng}
		 {\True $\overrightarrow{MP} = \overrightarrow{BC}$}
	\loigiai{ 
	\immini 
		 {
		Ta có $MN$ là đường trung bình của tam giác $ABC$ nên $MN = \dfrac{1}{2}BC$ và  $MN\parallel BC$.\\
		Do đó $MP = BC$ và $MP\parallel  BC$. Vậy
		\begin{itemchoice}
			\itemch 
			\itemch Vì $\left| \overrightarrow{MP} \right| = MP = BC =  \left| \overrightarrow{BC} \right|$.
			\itemch  Vì $\overrightarrow{MN}$ và $\overrightarrow{BC}$ cùng hướng.
			\itemch Vì $\left| \overrightarrow{MP} \right| = MP = BC = \left| \overrightarrow{BC} \right|$  và $\overrightarrow{MP}$, $\overrightarrow{BC}$ cùng hướng.
		\end{itemchoice}
		 }
		{ 
		\begin{tikzpicture}[line join=round,line cap = round,>=stealth,scale=.8,font=\footnotesize,transform shape]
				\foreach \x/\y/\z/\g in 
					 {
					0/0/B/-135,2/4/A/90,6/0/C/-45,1/2/M/135,4/2/N/60,7/2/P/45 
				 }
				\draw[fill=black] (\x,\y)  circle (1pt) coordinate (\z)
					 ($(\z)+(\g:3.5mm)$) node{$\z$};
				 \draw (B)--(A)--(C);
				\draw[->]  (M)--(P);
				\draw[->]  (M)--(N);
				 \draw[->] (B)--(C);
			\end{tikzpicture}
		} 
	} 
\end{ex}

% nguồn: MaHoaNH_DungSai_Lop10
\begin{ex}%[0H5H1-3]%[ma:7d723eb6]
	Cho tứ giác $ABCD$. Gọi $E$, $F$, $G$,  $H$ theo thứ tự là trung điểm của $BC$, $BD$, $AD$, $AC$.
	\choiceTF 
		{\True $EF$ là  đường trung bình của các tam giác $BCD$}
		 {$GH=\dfrac{1}{3} CD$}
		{\True $EFGH$ là hình bình  hành}
		{$\overrightarrow{HG}=\overrightarrow{FE}$} 
	\loigiai{ 
		\begin{center}
			\begin{tikzpicture}[declare function={r=3;}]
					\path 
						 (120:r) coordinate (A)
						 (45:r) coordinate (B)
						(-30:r) coordinate (C) 
						(210:r) coordinate (D) 
						 ($(B)!.5!(C)$) coordinate (E)
						($(B)!.5!(D)$) coordinate  (F)
						($(A)!.5!(D)$)  coordinate (G)
						($(A)!.5!(C)$) coordinate (H); 
					\draw (A)--(B)--(C)--(D)--cycle  (A)--(C)
						(B)--(D) 
						(E)--(F)--(G)--(H)--cycle; 
					\foreach  \t/\g in {A/120,B/45,C/-30,D/210,E/0,F/-60,G/180,H/60}{
						\draw[fill=black] (\t) circle (1pt) node[shift={(\g:9pt)},font=\scriptsize]{$  \t $};
					} 
				\end{tikzpicture}
		\end{center}
		\begin{itemchoice}
			\itemch Ta có $EF$ là đường  trung bình của tam giác $BCD$.\\
			Suy ra $EF\parallel CD$  và $EF=\dfrac{1}{2} CD$. \hfill (1)
			\itemch Ta có $GH$ là đường trung bình của các tam giác  $ACD$ nên
			$GH\parallel CD$ và $GH=\dfrac{1}{2}  CD$. \hfill (2)
			\itemch Từ  (1) và (2) suy ra $EFGH$ là hình bình hành.
			\itemch  $\overrightarrow{HG}=\overrightarrow{EF}$.
		\end{itemchoice}
	 }
\end{ex}


%%%%%%%%%%%%%%%% TRẢ LỜI NGẮN %%%%%%%%%%%%%%%%

% nguồn: MaHoaNH_TraLoiNgan_Lop10
\begin{bt}%[0H5H1-3]%[ma:ebc5291c]
	Cho lục giác đều $ABCDEF$  tâm $O$. Có bao nhiêu vectơ bằng vectơ $\overrightarrow{OD}$ có điểm đầu và điểm cuối là các điểm trong hình (không tính $\overrightarrow{OD}$)?
	 \shortans{3}
	 \loigiai{
	\immini 
		 {
		Do  $ABCDEF$ là lục giác đều nên $AO = OD = BC = FE$ và $BC\parallel AD\parallel FE$. \\
		Vậy $\overrightarrow{BC} = \overrightarrow{AO} = \overrightarrow{FE} = \overrightarrow{OD}$. 
		} 
		 {
		\begin{tikzpicture}[line join=round,line cap = round,>=stealth,scale=.8,font=\footnotesize,transform shape]
				 \foreach \x/\z/\g in
					{ 
					A/180/180,B/120/135,C/60/45,D/0/0,E/-60/-45,F/-120/-135 
				} 
				\draw[fill=black] ($(0,0)+(\z:3)$) circle (1pt) coordinate  (\x)
					($(\x)+(\g:3.5mm)$) node{$\x$}; 
				\draw[fill = black] (0,0)  circle (1pt) node[below]{$O$};
				\draw 
					 (A)--(B)--(C)--(D)--(E)--(F)--(A)
					(A)--(D) 
					(B)--(E) 
					 (C)--(F);
			\end{tikzpicture}
		} 
	} 
\end{bt}

% nguồn: MaHoaNH_TraLoiNgan_Lop10
\begin{bt}%[0H5V1-3]%[ma:b508da56]
	Cho $4$ điểm phân biệt. Có nhiều nhất bao nhiêu  cặp vectơ khác vectơ-không và bằng nhau được tạo thành từ $4$ điểm đã cho?
	 \shortans{8}
	\loigiai{ 
		\begin{center}
			\begin{tikzpicture}[line join=round,line cap = round,>=stealth,scale=.8,font=\footnotesize,transform shape]
					\foreach \x/\y/\z/\g  in
						{ 
						0/0/A/-90,2/0/B/-90,4/0/C/-90,6/0/D/-90 
					 }
					\draw[fill=black] (\x,\y) circle (1pt)  coordinate (\z)
						 ($(\z)+(\g:3.5mm)$) node{$\z$};
					\draw (A)--(B)--(C)--(D); 
				\end{tikzpicture}
		\end{center}
		Gọi $4$  điểm phân biệt là $A$, $B$, $C$, $D$.\\
		Để tạo ra nhiều nhất cặp vectơ khác vectơ-không và bằng nhau từ $4$ điểm  phân biệt thì ta xếp $4$ điểm $A$, $B$, $C$, $D$ trên một đường thẳng thoả mãn $AB = BC = CD$.\\
		Khi đó ta tạo được  $8$ cặp vectơ bằng nhau là $\overrightarrow{AB} = \overrightarrow{BC}$, $\overrightarrow{AB} = \overrightarrow{CD}$, $\overrightarrow{BC} = \overrightarrow{CD}$, $\overrightarrow{AC} = \overrightarrow{BD}$ và các cặp vectơ đối của mỗi cặp.
	 }
\end{bt}


%%%%%%%%%%%%%%%% TỰ LUẬN %%%%%%%%%%%%%%%%

% nguồn: MaHoaNH_TuLuan_Lop10
\begin{bt}%[0H5H1-3]%[ma:b78026fa]
	Cho tam giác $ABC$, trung tuyến $AM$. Trên cạnh $AC$ lấy điểm $N$ sao cho $AC=3AN$. Gọi $I$ là giao điểm của $AM$ và $BN$. Chứng  minh $\overrightarrow{AI}=\overrightarrow{IM}$.
	 \loigiai{
		\begin{center}
			\begin{tikzpicture}[>=stealth,line join=round,line cap=round,font=\footnotesize,scale=0.8]
					\path 
						(1.5,4) coordinate (A) 
						 (0,0) coordinate (B)
						 (5,0) coordinate (C)
						 ($(B)!0.5!(C)$) coordinate (M)
						($(A)!0.5!(M)$)  coordinate (I)
						(intersection of B--I and A--C) coordinate (N) 
						 ($(N)!0.5!(C)$) coordinate (E);
					\draw (A)--(B)--(C)--cycle; 
					\draw  (A)--(M)
						(B)--(N) 
						(M)--(E); 
					\foreach \x/\g in  {A/90,B/200,C/30,I/140,N/30,M/-90,E/45} \fill (\x) circle (1pt) +(\g:3mm) node[scale=.8] {$\x$};
				\end{tikzpicture}
		\end{center}
		Gọi $E$ là trung  điểm $NC$. Ta suy ra $AN=NE=EC=\dfrac{AC}{3}$.\\
		Xét tam giác $BNC$ có $ME$ là  đường trung bình, suy ra $ME\parallel BN$.\\
		Xét tam giác  $AME$ có
		$$\heva{&N\text{ là trung điểm } AE\\& NI \parallel ME.}$$ 
		Suy ra $I$ là trung điểm $AM$.\\ 
		Ta  có
		$$\heva{&AI=IM\\&\overrightarrow{AI}\text{ và } \overrightarrow{IM}\text{ cùng  hướng}.}$$
		Suy ra $\overrightarrow{AI}=\overrightarrow{IM}$. 
	 }
\end{bt}
